\chapter{Considerações Finais}

O presente estudo comparativo demonstrou que a interface $\beta$‑catenina/BCL9 apresenta
maior estabilidade e afinidade do que os complexos $\beta$‑catenina/TCF4 e $\beta$‑catenina/LEF1,
reforçando seu potencial como alvo terapêutico prioritário no câncer colorretal. A
abordagem in silico permitiu caracterizar detalhadamente os perfis de interação, sugerindo
que o desenvolvimento de peptídeos miméticos ou inibidores específicos dessa interface
pode constituir uma estratégia promissora para a supressão da via Wnt. Futuras
investigações devem incluir simulações dinâmicas e triagem de candidatos farmacológicos
para validar a viabilidade clínica desses achados.

As mutações encontradas para o vírus HPV-16, especialmente a H85Y comprometem a integridade estrutural do complexo E6-p53. Os picos de RMSD na dinâmica molecular foram maiores para wild-type, o que induziu maiores alterações conformacionais e inibição pela p53. Dessa forma, os resultados fornecem uma base molecular para a hipótese de que pacientes infectados com esta variante possam apresentar desfechos clínicos distintos. 


